
\documentclass{article} 
\usepackage{graphicx} % Required for inserting images 
\usepackage[a4paper, left=2cm, right=2cm, top=3cm, bottom=3cm]{geometry} 
\usepackage{amsmath} %usa questo per scrivere le formule matematiche 
\usepackage{parskip} %Questo pack elimina l'indentazione fastidiosa  
\usepackage{amssymb} 
\usepackage{amsthm} 

% --- NUOVO PACCHETTO PER I BOX (Sostituisce mdframed/framed) ---
\usepackage[most]{tcolorbox}
\tcbuselibrary{theorems} % Permette di gestire i teoremi direttamente dentro i box

% --- DEFINIZIONE DEI COLORI ---
\definecolor{winered}{rgb}{0.6,0,0}      % Rosso scuro per Teoremi e Proposizioni
\definecolor{navyblue}{rgb}{0,0,0.5}     % Blu scuro per Definizioni
\definecolor{darkgray}{rgb}{0.3,0.3,0.3} % Grigio per Note ed Esempi

% --- CREAZIONE DEGLI AMBIENTI BOXATI ---
% Questa sintassi crea direttamente il box colorato e gestisce la numerazione in modo sicuro.
% I nomi degli ambienti (theo, prop, defi, ecc.) sono IDENTICI ai tuoi originali.

% 1. Teoremi (Rosso scuro, numerazione progressiva standard)
\newtcbtheorem[number within=section]{theo}{Theorem}{
    colback=winered!5, colframe=winered, fonttitle=\bfseries,
    colupper=black, coltitle=black,
    sharp corners, boxrule=0pt, leftrule=3pt,
    enhanced, attach title to upper
}{th}

% 2. Proposizioni (Stesso stile di Theorem)
\newtcbtheorem[use counter from=theo]{prop}{Proposition}{
    colback=winered!5, colframe=winered, fonttitle=\bfseries,
    colupper=black, coltitle=black,
    sharp corners, boxrule=0pt, leftrule=3pt,
    enhanced, attach title to upper
}{pr}

% 3. Corollari (Numerati in base al teorema corrente)
\newtcbtheorem[number within=theo]{crlr}{Corollario}{
    colback=winered!3, colframe=winered, fonttitle=\bfseries,
    colupper=black, coltitle=black,
    sharp corners, boxrule=0pt, leftrule=3pt,
    enhanced, attach title to upper
}{co}

% 4. Definizioni (Blu scuro)
\newtcbtheorem[number within=section]{defi}{Definition}{
    colback=navyblue!5, colframe=navyblue, fonttitle=\bfseries,
    colupper=black, coltitle=black,
    sharp corners, boxrule=0pt, leftrule=3pt,
    enhanced, attach title to upper
}{def}

%\newtcbtheorem{proof}{proof}{}{proof}

% 5. Note (Grigio, senza numerazione)
\newtcolorbox{note}{
    colback=darkgray!5, colframe=darkgray,
    sharp corners, boxrule=0pt, leftrule=3pt,
    enhanced, before upper={\textbf{Note. }}
}

% 6. Esempi (Grigio, senza sfondo, solo linea a sinistra, senza numerazione)
\newtcolorbox{xpl}{
    colback=white, colframe=darkgray,
    sharp corners, boxrule=0pt, leftrule=2pt,
    enhanced, before upper={\textbf{Example. }}
}

\newenvironment{pop}{}{} 


\begin{document}

\section{Introduction} 
  In this article I'll write about electromagnetics.

  The iter will be:
\begin{enumerate}   
  \item Electric Fields,   
  \item Capacitors,   
  \item Dielectrics,    
  \item Current,   
  \item Resistors,   
  \item Magnetic Fields,   
  \item Electromagnetics.  
\end{enumerate}
\newpage{}

\section{Electric Fields}
We'll start by defining what is the force that applies to two charges $q_1, q_2$.
\begin{defi}{Force}{e-force}
  
  A \textbf{Force} between two charges $q_1, q_2$ is defined by the relation:
  $$ F = k_0 \frac{q_1 q_2}{r^2} $$
  This definition is in module, so $F = |\vec F|$. The force vector is defined by the rule:
  $$ \vec{F}_{1,2} = k_0 \frac{q_1 q_2}{(\vec{r}_{2,1})^2}\hat{r}_{2,1} $$
  This is the force that the charge $q_2$ applies to the charge $q_1$. (the vector $\vec{r}_{2,1}$ is the vector from $q_2 \text{ to } q_1$)\\
  It's also to remark that:
  $$ \vec{F}_{1,2} = -  \vec{F}_{2,1}$$
\end{defi}
\begin{defi}{Electric Field}{e-field}

  The \textbf{Field} of a charge $q$ is defined by the limit:
  $$ \vec{E} = \lim_{q \rightarrow 0} \frac{\vec{F}}{q} $$
  That definition is equivalent to the more practical: 
  $$ \vec{E}(\vec{r}) = k_0 \frac{q}{\vec{r}^2} \hat{r} $$
\end{defi}
\begin{note}
  Let's suppose that there are two charges $q_1$ and $q_2$, and we choose a random point in the space $p$ diffent 
  from the two charges. The field in that point will be: 
  $$ \Vec{E}_{Tot}(p) = \vec{E}_{q_1}(p) + \vec{E}_{q_2}(p) $$
\end{note}
\begin{note}
  We use the SI, so the force has [N] as unit, the charge has [C]-oulomb
 and the field has [N/C] or [$V/m$].
\end{note}
\begin{prop}{Sovrapposition}{label}

  The field in a random point p is given by:
  \begin{enumerate}
    \item For a discrete system of charges:
      $$ \vec{E}(p) = \sum^N_{i=1}\vec{E}_{q_i}(\vec{r_i}) $$
    \item For a continuos system of charges:
      $$ \vec{E}(p) = \int_Q \vec{E}(\vec{r})dQ $$
  \end{enumerate}
\end{prop}
\begin{proof}
  Just using the note seen before.

\end{proof}
\newpage{}
\begin{xpl}{(Field generated by two charges)}

  Let's put two charges, $q_1$ and $q_2$, in a plane. \\ 
  $q_1 = (0, y), q_1 = q, \quad q_2 = (0, -y), q_2 = -q, \quad y>0, q>0 [C]$.

  As you can see, I'll write the position of a charge in the same notation as the value of a charge (when it's clear by the context).
  
  In this example I'll find the value of the field in the point $p=(x,0), \quad x>0.$

  Let's see the field of the charge $q_1$ on the point $p$.  
  $$\vec{E}_{q_1}(p)=k_0 \frac{q_1}{\overrightarrow{r_{q_1}}^2}\hat{r}_{q_1} = 
  k_0\frac{q_1}{{x^2 +(-y)^2}}(\frac{x}{\sqrt{x^2+(-y)^2}},\frac{-y}{\sqrt{x^2+(-y)^2}}) $$
  So we can express the field in the coordinates $x,y$: 
  $$ \overrightarrow{E_{q_1}}(p) = k_0 (\frac{q}{(x^2 + y^2)^{3/2}}x,\frac{-q}{(x^2 + y^2)^{3/2}}y) $$
  The field in $p$ of the charge $q_2$ is:
  $$ \overrightarrow{E_{q_2}}(p) = k_0 (\frac{-q}{(x^2 + y^2)^{3/2}}x,\frac{-q}{(x^2 + y^2)^{3/2}}y) $$
  So the total field in $p$ is: 
  $$ \overrightarrow{E_{Tot}}(p) = \overrightarrow{E_{q_2}}(p) + \overrightarrow{E_{q_1}}(p) 
  = (0,k_0\frac{-2q}{(x^2 + y^2)^{3/2}}y) $$
\end{xpl}
\newpage{}  
\begin{xpl}{Field generated by a charged string}

  Let $s$ be an infinite string. I want to describe the field in this system.\\ 
  Let $s$ be in the $z$ axis. There's an important symmetrical situation.\\ 
  If the coordinates (x,y) are fixed, the module of the field is independent from z.\\ 
  furthermore, the module of the field is the same in the circumference of radius r.\\ 
  That's why it's a good idea to study the field in the coordinates (z,d) where d is the distance from the string.
  
  For every $dz$ there's a charge equals to $\lambda dz$, where $\lambda$ is the charge density of the string.\\ 
  Let's fix a point $p=(0,d)$ in the coordinates $(z,d)$.
  $$ d \overrightarrow{E_{dz}}(p) = k_0\frac{\lambda dz}{\overrightarrow{r_{dz}}^3}\overrightarrow{r_{dz}} $$
  The simmetry of the system tells that the z-contribute of every point z is same in module 
  but different sign of the pont -z. For this reason it's resonable to evaluate 
  only the field on the coordinate orizzontaly.
  $$ d\overrightarrow{E_d}(p) =  k_0\frac{\lambda dz}{(z^2 + d^2)^2}\frac{d}{\sqrt{z^2 + d^2}}$$
  Now I can evaluate the field integrating on the z axis.
  \begin{align*}
    \overrightarrow{E_d}(p) &=\int^{+\infty}_{-\infty}d\overrightarrow{E_d}(p)dz \\ 
                            &=\int^{+\infty}_{-\infty}k_0\frac{\lambda d}{(z^2 + d^2)^{3/2}}dz 
  \end{align*}
  \begin{note}
    An elegant algebraic method consists in factoring out $z^2$ from the denominator:
    \begin{align*}
      \int \frac{dz}{(z^2 + d^2)^{3/2}} &= \int \frac{dz}{z^3 \left(1 + \frac{d^2}{z^2}\right)^{3/2}}
    \end{align*}
    By applying the algebraic substitution $u = 1 + \frac{d^2}{z^2}$, we get $du = -\frac{2d^2}{z^3}dz$, which leads to:
    \begin{align*}
      \int \frac{dz}{(z^2 + d^2)^{3/2}} &= -\frac{1}{2d^2} \int u^{-3/2} du \\
                                   &= -\frac{1}{2d^2} \left( \frac{u^{-1/2}}{-1/2} \right) = \frac{1}{d^2\sqrt{u}} + C \\
                                   &= \frac{z}{d^2\sqrt{z^2 + d^2}} + C
    \end{align*}
  \end{note}
  At this point we have a result
    \begin{align*}
      \overrightarrow{E_d}(p) &= [k_0 \lambda d \frac{z}{d^2\sqrt{z^2 + d^2}}]^{+\infty}_{-\infty}  \\ 
                              &= k_0 \lambda d \frac{2}{d^2} \\
                              &= \frac{2 k_0 \lambda }{d} \\ 
                              &= \frac{\lambda}{2 \pi \epsilon_0 d}
  \end{align*} 
\end{xpl}
\newpage{}
\begin{theo}{Gauss Theorem}{gauss}

  The flux of $\overrightarrow{E}$ in the empty space of a closed surface $S$ is given by the sum of 
  the charges inside $S$ divided by $\epsilon_0 $.

  $$ \Phi (\overrightarrow{E}) = \frac{Q^{inside}_{Tot}}{\epsilon_0 } $$
  Using the definitions it becomes: 
  $$ \int_S \overrightarrow{E} \hat{n} d\sigma = \int_{V} \frac{\rho(x)}{\epsilon_0 } d\tau $$ 
\end{theo}
\begin{proof}
  Let's take the elementar part of the flux of a charge inside $S$. 
  \begin{align*}
      d \Phi \overrightarrow{E} &= \overrightarrow{E} \hat{n} d\sigma \\
                                &= \overrightarrow{E} d\vec{\sigma} \\
                                &= E cos \theta d\sigma \\ 
                                &= E d\sigma_n \\ 
                                &= \frac{Q}{4\pi\epsilon_0}\frac{d\sigma_n}{r^2}
\end{align*}
The element of surface is equal to $d\sigma_n = r^2 sin \theta d\theta d\phi $. So the element of flux is 
$$ d\Phi\overrightarrow{E} =  \frac{Q}{4\pi\epsilon_0} sin\theta d\theta d\phi  $$

So the flux of this charge $Q$ will be 
\begin{align*}
  \Phi(\overrightarrow{E}) &= \int_S d\Phi (\overrightarrow{E}) \\ 
                        &= \frac{Q}{4\pi\epsilon_0}\int^{2\pi}_{0}d\phi \int^{\pi}_{0}sin\theta d\theta \\ 
                        &= \frac{Q}{\epsilon_0}
\end{align*}
For the principle of Sovrapposition the flux is the sum of the flux of the internal charges. \\ 
The contribute of the external charges is null due to the fact that the flux that enters the surface 
is the same of the flux that exits.

\end{proof}


\end{document}

